\begin{frame}{Длинная таблица: перечень пакетов\orient}
  \vspace{0.2em}
  {\footnotesize\color{acGray}Пример компактной таблицы с {\ttfamily\small} или
  {\ttfamily footnotesize}.}
  \vspace{0.5em}

  \footnotesize
  \begin{tabular}{@{}p{0.22\textwidth}p{0.74\textwidth}@{}}
    \toprule
    \textbf{Категория} & \textbf{Пакеты / компоненты}\\
    \midrule
    Основные &
      \ttfamily component-a, component-b, component-c, component-d \\
    \addlinespace[0.3em]
    Обработка данных &
      \ttfamily lib-1, lib-2, lib-3, lib-4 \\
    \addlinespace[0.3em]
    Визуализация &
      \ttfamily chart-a, chart-b, chart-c \\
    \addlinespace[0.3em]
    Утилиты &
      \ttfamily util-1, util-2, util-3 \\
    \bottomrule
  \end{tabular}
\end{frame}

\begin{frame}{Два блока: условия и практика}
  \vspace{0.3em}
  \begin{columns}[T]
    \begin{column}{0.49\textwidth}
      \begin{block}{Обязательные условия}
        \begin{itemize}\setlength\itemsep{0.3em}
          \item \textbf{прямой доступ} к настройкам оборудования;
          \item пропускная способность канала --- \textbf{не менее 200 Мбит/с}.
        \end{itemize}
      \end{block}
    \end{column}
    \begin{column}{0.49\textwidth}
      \begin{block}{На практике}
        \begin{itemize}\setlength\itemsep{0.3em}
          \item выделенный сегмент сети;
          \item ограниченный доступ в Интернет;
          \item тестовый прогон до основного события.
        \end{itemize}
      \end{block}
    \end{column}
  \end{columns}
\end{frame}
